% Adaptación de ensamblado: espacio de nombres para destinos PDF.
\renewcommand{\theHchapter}{fisica.\arabic{chapter}}
% !TeX root = ../main.tex
%% ============================================================================
%%  04_arquitectura_fisica/contenido.tex
%%
%%  Subdocumento 4.2 del Formulario T-7 --- Arquitectura física y de despliegue.
%%  Se monta desde main.tex con  \subdocumento{04_arquitectura_fisica}.
%%
%%  La carpeta es autocontenida:
%%    partes/    las secciones, una por archivo, traídas con \parte{}
%%    figuras/   el modelo de emplazamiento y los planos del recinto
%%
%%  Ninguna ruta de aquí sube con ../ ni parte de la raíz del repositorio, de
%%  modo que mover la carpeta no rompe nada.
%%
%%  NUMERACIÓN. Cada pieza es un capítulo de la clase: capítulos 1..15,
%%  secciones x.y y subsecciones x.y.z. El título de la parte no numera: es
%%  un rótulo de apertura, de modo que ninguna sección del documento arrastra
%%  el prefijo "1." de una sección contenedora. La losa (a)-(n) fue eliminada;
%%  los títulos de las piezas nombran su contenido y las referencias internas
%%  citan por nombre de sección, no por símbolo §.
%% ============================================================================

\renewcommand{\thechapter}{4.\arabic{chapter}}
\setcounter{chapter}{1}   % 4.1 es la arquitectura lógica; la física abre como 4.2

% Títulos de apartado como "4.2 Arquitectura física", en una línea y sin el
% rótulo "Capítulo" que pone lafrox.cls. Mismo tamaño y color del título de
% la clase. La antigua "Visión general" (partes/01_vision_general.tex) quedó
% como introducción de 4.2 en partes/02_0_arquitectura_fisica.tex.
\titleformat{\chapter}[block]
  {\normalfont\raggedright\intersemibold\fontsize{25pt}{28pt}\selectfont\color{lafroxFuerte}}
  {\thechapter}{0.45em}{}

% --- 4.2 Arquitectura física (índice obligatorio, aclaraciones sección 11) ---
%  02_0 abre el apartado con número "4.2". El índice obligatorio fija 4.2.1 =
%  "Especificaciones Implementos a proveer (Hardware y Software)"; los
%  contenidos de 4.2 (emplazamiento, servicios, despliegue, conexiones y
%  dimensionamiento) siguen como 4.2.2 a 4.2.6 (aclaraciones, sección 2).
%  Las etiquetas de sección no llevan número, así las referencias no dependen
%  del orden.
\parte{02_0_arquitectura_fisica}
\parte{04_c_implementos}        % 4.2.1 (título obligatorio)
\parte{02_a_emplazamiento}      % 4.2.2
\parte{02_b_servicios_nube}     % 4.2.3
\parte{11_j_despliegue}         % 4.2.4
\parte{02_c_conexiones}         % 4.2.5
\parte{12_k_dimensionamiento}   % 4.2.6
\let\LFXcuerpotabla\LFXcuerpotablaorig   % tablas sin guiones solo en 4.2
% --- fin del bloque 4.2 ---

% --- 4.3 Data center: primario (4.3.1) y secundario (4.3.2) ---
\parte{../4.3_centros_de_datos/05_d_sitio_principal}    % 4.3 y 4.3.1 data center primario
\parte{../4.3_centros_de_datos/06_e_sitio_secundario}   % 4.3.2 data center secundario
% 07_f_niveles_servicio: cubierto en el sitio principal (99,95 % por componente) y en alta disponibilidad (>= 99,9 % transacción crítica)
% 08_g_operacion_desconectada: absorbido en 4.2.1 (autonomía sin enlace, tab:t29)
% La arquitectura de integración y la de seguridad pertenecen a la arquitectura lógica (apartado 4.1).
% 13_l_capa_analitica: eliminado. La capa analítica pertenece a arquitectura lógica (Subdocumento 4.1); sus componentes físicos están en la síntesis de emplazamiento (tab:t21) y en los apartados de componentes de nube pura (N-04, N-06, N-08, N-10)
% 03_b_tecnologias: retirado; las tecnologías de software se especifican en el apartado 4.1.1 (arquitectura lógica)
\parte{14_m_decisiones_adr}
\parte{17_referencias}
% 15_n_trazabilidad: la trazabilidad código por código se responde en el Formulario T-12 (Subdocumento 3); las referencias cruzadas están dentro de cada apartado

% Formulario T-11, provisoriamente al final del subdocumento para tenerlo a
% mano. Para la entrega va como archivo propio (LAFROX-Formulario-T-11.tex,
% aclaraciones, sección 1): basta comentar la línea \parte de abajo. Ambos
% usan el mismo cuerpo, partes/15_anexo_t11.tex.
\renewcommand{\LFXcuerpotabla}{%   % el anexo no corta en guiones explícitos (RT-02.11, cross-docking)
  \LFXcuerpotablaorig\exhyphenpenalty=10000\relax}
\parte{../../formularios/T11/15_anexo_t11}
% Anexo 4.B: memoria de cálculo del dimensionamiento, después del T-11.
\parte{../../anexos/fisica/16_anexo_4b_memoria_calculo}
